% =============================================================================
% APPENDIX CP: DOMAIN-CPS: Complex Problem Solving - 10C Integration
% =============================================================================
% Category: DOMAIN
% Module: BCM2_DOMAIN_CPS
% Version: 1.0 (January 2026)
% Status: Final
% Language: English
%
% This appendix integrates the Complex Problem Solving (CPS) research tradition
% (Dörner, Funke, Greiff) into the EBF 10C framework. CPS is modeled not as a
% separate theory but as a specific CONFIGURATION of existing CORE dimensions.
%
% Session: EBF-S-2026-01-28-EDU-001
% Model: MOD-CPS-001
% =============================================================================

% =============================================================================
% CROSS-REFERENCE STRUCTURE
% =============================================================================
%
% CHAPTER LINKAGE:
% - Primary Chapter: Chapter 11 (CORE-AWARE: The Awareness Filter)
% - Secondary Chapters:
%   - Chapter 9 (CORE-WHEN: Context Framework)
%   - Chapter 4 (CORE-HOW: Complementarity)
%   - Chapter 15 (CORE-HIERARCHY: Decision Hierarchy)
%
% DEPENDENCIES (what this appendix REQUIRES):
% - Appendix AU (CORE-AWARE): Awareness function A(t*)
% - Appendix V (CORE-WHEN): Context dimensions Ψ
% - Appendix B (CORE-HOW): Complementarity γ
% - Appendix HI (CORE-HIERARCHY): Decision levels L0-L3
% - Appendix C (CORE-WHAT): Utility dimensions (U_D)
% - Appendix AV (CORE-READY): Willingness WAX, threshold θ
%
% DEPENDENTS (what REQUIRES this appendix):
% - Domain applications in Education, Management, Training
% - Appendix HHH (METHOD-TOOLKIT): Intervention design for CPS contexts
%
% =============================================================================

\chapter{Complex Problem Solving: 10C Integration}
\label{app:domain-cps}

% -----------------------------------------------------------------------------
% CHAPTER LINKAGE BOX
% -----------------------------------------------------------------------------

\begin{tcolorbox}[colback=purple!5!white, colframe=purple!75!black,
    title=Chapter Linkage]

\textbf{Primary Chapter:} Chapter 11: The Awareness Filter (CORE-AWARE)
\begin{itemize}[nosep]
\item Section 11.3: Bounded Awareness $\leftrightarrow$ This Appendix Section \ref{sec:cps-aware}
\item Section 11.5: Cognitive Load $\leftrightarrow$ This Appendix Section \ref{sec:cps-parameters}
\end{itemize}

\textbf{Secondary Chapters:}
\begin{itemize}[nosep]
\item Chapter 9: Context Framework --- Ψ dimensions for CPS (intransparency, dynamics)
\item Chapter 4: Complementarity --- γ connectivity between problem variables
\item Chapter 15: Decision Hierarchy --- L0-L3 levels in strategic problem solving
\end{itemize}

\textbf{Reading Guide:}
\begin{itemize}[nosep]
\item For conceptual overview: Read Chapter 11 (AWARE) first
\item For context dimensions: See Appendix V (CORE-WHEN)
\item For formal Awareness function: See Appendix AU
\end{itemize}

\end{tcolorbox}

% -----------------------------------------------------------------------------
% ABSTRACT
% -----------------------------------------------------------------------------

\begin{abstract}
Complex Problem Solving (CPS) is a well-established research paradigm originating from Dörner's work on the ``Logic of Failure'' and operationalized through Funke's MicroDYN approach. This appendix demonstrates that CPS is not a separate theory in the EBF but rather a specific \textbf{configuration of 6 existing CORE dimensions}: AWARE, HOW, WHEN, READY, STAGE, and WHAT. We provide the core formula $U_{\text{CPS}}(t) = A(\Psi_{\text{load}}) \times (1-\eta) \times [U_D + \sum \gamma_{ij} \times (\frac{dX_i}{dt} \times \frac{dX_j}{dt})]$, validate parameters from the CPS literature, and derive intervention priorities based on sensitivity analysis showing intransparency (η) as the primary lever (42\%).
\end{abstract}

% -----------------------------------------------------------------------------
% QUICK REFERENCE BOX
% -----------------------------------------------------------------------------

\begin{tcolorbox}[colback=blue!5!white, colframe=blue!75!black,
    title=Quick Reference: CPS in 10C]

\textbf{Core Formula:}
\begin{equation}
U_{\text{CPS}}(t) = A(\Psi_{\text{load}}) \times (1-\eta) \times \left[U_D + \sum_{i,j} \gamma_{ij} \times \frac{dX_i}{dt} \times \frac{dX_j}{dt}\right]
\label{eq:cps-utility}
\end{equation}

\textbf{Key Parameters:}
\begin{itemize}[nosep]
\item $A(\Psi_{\text{load}})$: Awareness under cognitive load (CORE-AWARE)
\item $\eta \in [0,1]$: Intransparency of problem system (CORE-WHEN)
\item $\gamma_{ij}$: Connectivity between variables (CORE-HOW)
\item $dX/dt$: System dynamics (CORE-WHEN)
\item $U_D$: Development utility / competence (CORE-WHAT)
\item $\theta$: Action threshold (CORE-READY)
\end{itemize}

\textbf{Action Condition:}
\begin{equation}
\text{Action} \iff WAX(U_{\text{CPS}}, \delta) \geq \theta(\Psi_{\text{uncertainty}})
\end{equation}

\textbf{Sensitivity (Intervention Priority):}
η↓ (42\%) → γ→ (28\%) → A↑ (17\%) → δ↓ (10\%) → θ (3\%)

\end{tcolorbox}

% -----------------------------------------------------------------------------
% APPENDIX SCOPE BOX
% -----------------------------------------------------------------------------

\begin{tcolorbox}[colback=green!5!white, colframe=green!75!black,
    title=Appendix Scope]

\textbf{Objective:} Integrate the CPS research tradition into the EBF 10C framework by showing that CPS emerges as a configuration of existing CORE dimensions rather than requiring new theoretical primitives.

\textbf{In-Scope:}
\begin{itemize}[nosep]
\item Mapping Dörner's 5 qualities to 10C dimensions
\item Formal specification of CPS utility function
\item Parameter calibration from CPS literature (Funke, Greiff, Sweller)
\item Sensitivity analysis and intervention priorities
\item Comparison with PISA 2015 CPS framework
\end{itemize}

\textbf{Out-of-Scope:}
\begin{itemize}[nosep]
\item Detailed psychometric properties of MicroDYN (see \cite{PAP-greiff2013cps})
\item Collaborative Problem Solving (see PISA 2015 for social dimensions)
\item Domain-specific CPS applications (see worked examples for illustrations)
\end{itemize}

\textbf{Deliverables:}
\begin{itemize}[nosep]
\item 5 Axioms (CPS-1 to CPS-5) formalizing the 10C-CPS mapping
\item Parameter table with literature sources and confidence intervals
\item Intervention priority ranking based on sensitivity analysis
\end{itemize}

\end{tcolorbox}


% =============================================================================
% PART 2: CORE CONTENT
% =============================================================================

\section{Motivation: The Fundamental Question}
\label{sec:cps-motivation}

Complex Problem Solving (CPS) addresses situations where individuals face problems characterized by \textbf{complexity}, \textbf{connectivity}, \textbf{dynamics}, \textbf{intransparency}, and \textbf{polytely} (multiple goals). The fundamental question is:

\begin{tcolorbox}[colback=yellow!10!white, colframe=orange!75!black,
    title=Central Question]
How can the established CPS research paradigm (Dörner, Funke, Greiff) be integrated into the EBF 10C framework, and what does this integration reveal about intervention priorities?
\end{tcolorbox}

The key insight is that CPS does not require new theoretical primitives. Instead, it emerges as a \textbf{specific configuration} of existing CORE dimensions where:
\begin{itemize}
\item Awareness is bounded by cognitive load (AWARE)
\item Problem variables are interconnected (HOW: γ > 0)
\item The system exhibits dynamics and intransparency (WHEN: Ψ)
\item Action thresholds are elevated due to uncertainty (READY)
\item Problem-solving proceeds through stages (STAGE)
\item Competence motivation drives engagement (WHAT: $U_D$)
\end{itemize}


\section{Historical Background: From Lohhausen to MicroDYN}
\label{sec:cps-history}

The CPS research tradition began with Dörner's seminal work:

\begin{enumerate}
\item \textbf{Lohhausen Study (1983)}: Participants acted as mayor of a simulated town with 2000+ interconnected variables over 10 simulated years \citep{PAP-doerner1983lohhausen}.

\item \textbf{Logic of Failure (1989)}: Dörner identified systematic errors in complex situations: failure to consider side effects, tendency toward ``repair service behavior,'' and difficulty with exponential dynamics \citep{PAP-doerner1989logik}.

\item \textbf{MicroDYN Approach (2009-2017)}: Funke and Greiff developed minimal complex systems for standardized assessment, using linear equations with 3 input and 3 output variables \citep{PAP-greiff2017microdynamapproach}.

\item \textbf{PISA Integration (2012-2015)}: CPS became part of international large-scale assessments, with PISA 2015 assessing collaborative problem solving in 52 countries \citep{PAP-oecd2017pisa2015cps}.
\end{enumerate}


\section{Dörner's Five Qualities and 10C Mapping}
\label{sec:cps-five-qualities}

Dörner (1980, 1989) identified five defining qualities of complex problems. Table \ref{tab:cps-mapping} shows their mapping to 10C dimensions:

\begin{table}[htbp]
\centering
\caption{Mapping Dörner's 5 Qualities to 10C Dimensions}
\label{tab:cps-mapping}
\begin{tabular}{llll}
\toprule
\textbf{Dörner Quality} & \textbf{Definition} & \textbf{10C CORE} & \textbf{Parameter} \\
\midrule
Complexity & Many variables & AWARE & $A(\Psi_{\text{load}})$ \\
Connectivity & Variables influence each other & HOW & $\gamma_{ij} > 0$ \\
Dynamics & System changes autonomously & WHEN & $dX/dt$, $\Psi_T$ \\
Intransparency & Hidden states and relationships & WHEN & $\eta \in [0,1]$ \\
Polytely & Multiple, conflicting goals & WHAT & Multiple $U_d$ \\
\bottomrule
\end{tabular}
\end{table}


\section{The 10C-CPS Configuration}
\label{sec:cps-configuration}

CPS is formally defined as the following configuration of CORE dimensions:

\begin{tcolorbox}[colback=gray!10!white, colframe=gray!75!black,
    title=Definition: CPS as 10C Configuration]

\textbf{CPS} $:=$ \{
\begin{itemize}[nosep]
\item AWARE: $A(\Psi_{\text{load}}) < 1$ \hfill [bounded cognitive capacity]
\item HOW: $\gamma_{ij} > 0$ for connected variables \hfill [interconnected system]
\item WHEN: $\Psi_T$ dynamic, $\eta > 0$, $\delta > 0$ \hfill [dynamics, opacity, delay]
\item READY: $\theta(\Psi_{\text{uncertainty}}) \uparrow$ \hfill [elevated threshold]
\item STAGE: $S(t)$: Exploration → Hypothesis → Test → Solution
\item WHAT: $U_D > 0$ \hfill [competence motivation]
\item HIERARCHY: L0 → L1 → L2 decision levels
\}
\end{itemize}

\end{tcolorbox}

\subsection{AWARE: Bounded Cognitive Capacity}
\label{sec:cps-aware}

Following Cognitive Load Theory \citep{PAP-sweller1988cognitiveload}, awareness under load is:

\begin{equation}
A(\Psi_{\text{load}}) = A_{\max} \times \exp(-k \times \Psi_{\text{load}})
\label{eq:awareness-load}
\end{equation}

where:
\begin{itemize}
\item $A_{\max} = 0.85$ (maximum awareness capacity)
\item $k \approx 1.5$ (decay rate)
\item $\Psi_{\text{load}} \in [0,1]$ (cognitive load level)
\end{itemize}

\textbf{Implications:}
\begin{itemize}
\item At low load ($\Psi_{\text{load}} = 0.3$): $A = 0.54$
\item At high load ($\Psi_{\text{load}} = 0.7$): $A = 0.30$
\item At maximum load ($\Psi_{\text{load}} = 1.0$): $A = 0.19$
\end{itemize}


\subsection{HOW: Connectivity and Complementarity}
\label{sec:cps-how}

In complex problems, variables are interconnected. A change in variable $X_i$ affects variable $X_j$ through the complementarity parameter $\gamma_{ij}$:

\begin{equation}
\Delta X_j = \gamma_{ij} \times \Delta X_i + \epsilon
\end{equation}

The \textbf{average connectivity} in typical CPS scenarios:
\begin{equation}
\gamma_{\text{avg}} = 0.45 \pm 0.12 \quad \text{(Funke, 2001)}
\end{equation}


\subsection{WHEN: Intransparency, Dynamics, and Delay}
\label{sec:cps-when}

Three WHEN-dimensions are critical for CPS:

\textbf{1. Intransparency (η):}
\begin{equation}
A_{\text{eff}} = A(\Psi_{\text{load}}) \times (1 - \eta)
\end{equation}

where $\eta \in [0,1]$ captures:
\begin{itemize}
\item $\eta_{\text{start}}$: Initial state opacity (what is the current situation?)
\item $\eta_{\text{process}}$: Process opacity (what are the effects of actions?)
\end{itemize}

\textbf{2. Dynamics (dX/dt):}
\begin{equation}
\frac{dX_i}{dt} = f_i(X_1, \ldots, X_n, u_t)
\end{equation}

The system changes autonomously even without intervention ($u_t = 0$).

\textbf{3. Feedback Delay (δ):}
\begin{equation}
Y_{\text{observed}}(t) = Y_{\text{actual}}(t - \delta)
\end{equation}

Delays of $\delta = 3$--$8$ steps are typical in MicroDYN assessments.


% =============================================================================
% AXIOMS
% =============================================================================

\section{Axioms: Formal CPS-10C Integration}
\label{sec:cps-axioms}

\begin{tcolorbox}[colback=red!5!white, colframe=red!75!black,
    title=Axiom CPS-1: Configuration Principle]
Complex Problem Solving is not a separate theory but a \textbf{configuration} of existing 10C CORE dimensions. No new theoretical primitives are required.
\end{tcolorbox}

\begin{tcolorbox}[colback=red!5!white, colframe=red!75!black,
    title=Axiom CPS-2: Multiplicative Transparency]
Intransparency (η) enters the utility function \textbf{multiplicatively}, not additively, because complete opacity ($\eta = 1$) renders all other factors irrelevant:
\begin{equation}
U_{\text{CPS}} \propto (1 - \eta)
\end{equation}
This follows EXC-5 Whitelist criterion V1: If $\eta = 1$, then $U_{\text{CPS}} = 0$.
\end{tcolorbox}

\begin{tcolorbox}[colback=red!5!white, colframe=red!75!black,
    title=Axiom CPS-3: Awareness-Load Decay]
Awareness decays exponentially with cognitive load:
\begin{equation}
A(\Psi_{\text{load}}) = A_{\max} \times e^{-k \times \Psi_{\text{load}}}
\end{equation}
This functional form is justified by working memory capacity constraints \citep{PAP-sweller1988cognitiveload}.
\end{tcolorbox}

\begin{tcolorbox}[colback=red!5!white, colframe=red!75!black,
    title=Axiom CPS-4: Connectivity Interaction]
The effect of system dynamics on utility is modulated by connectivity:
\begin{equation}
\Delta U \propto \sum_{i,j} \gamma_{ij} \times \frac{dX_i}{dt} \times \frac{dX_j}{dt}
\end{equation}
Higher connectivity ($\gamma_{ij} > 0$) amplifies the impact of dynamic changes.
\end{tcolorbox}

\begin{tcolorbox}[colback=red!5!white, colframe=red!75!black,
    title=Axiom CPS-5: Threshold Elevation Under Uncertainty]
The action threshold θ increases with problem uncertainty:
\begin{equation}
\theta(\Psi_{\text{uncertainty}}) = \theta_0 + \kappa \times \Psi_{\text{uncertainty}}
\end{equation}
where $\theta_0 \approx 0.55$ (satisficing baseline, Simon 1955) and $\kappa > 0$.
\end{tcolorbox}


% =============================================================================
% KEY RESULTS
% =============================================================================

\section{Key Results and Propositions}
\label{sec:cps-results}

The 10C-CPS integration yields the following key results:

\begin{tcolorbox}[colback=blue!5!white, colframe=blue!75!black,
    title=Proposition CPS-R1: Configuration Sufficiency]
Complex Problem Solving can be \textbf{fully characterized} within the EBF 10C framework using 6 CORE dimensions (AWARE, HOW, WHEN, READY, STAGE, WHAT) without requiring additional theoretical primitives.
\end{tcolorbox}

\begin{tcolorbox}[colback=blue!5!white, colframe=blue!75!black,
    title=Proposition CPS-R2: Intransparency Dominance]
Among all CPS parameters, intransparency ($\eta$) has the \textbf{largest marginal effect} on CPS utility:
\begin{equation}
\frac{\partial U_{\text{CPS}}}{\partial \eta} > \frac{\partial U_{\text{CPS}}}{\partial \gamma} > \frac{\partial U_{\text{CPS}}}{\partial A} > \frac{\partial U_{\text{CPS}}}{\partial \delta}
\end{equation}
Implication: Interventions should prioritize transparency enhancement over other approaches.
\end{tcolorbox}

\begin{tcolorbox}[colback=blue!5!white, colframe=blue!75!black,
    title=Proposition CPS-R3: Multiplicative Structure]
The CPS utility function is \textbf{multiplicative} in transparency $(1-\eta)$ and awareness $A(\Psi_{\text{load}})$, implying that neither can compensate for the other. If $\eta = 1$ or $A = 0$, then $U_{\text{CPS}} = 0$ regardless of other parameters.
\end{tcolorbox}

\begin{tcolorbox}[colback=blue!5!white, colframe=blue!75!black,
    title=Proposition CPS-R4: EBF-PISA Complementarity]
The EBF 10C approach and the PISA 2015 CPS framework are \textbf{complementary}: EBF explains \textit{why} (mechanisms), PISA measures \textit{what} (competencies). Integration enables mechanism-based interpretation of PISA scores.
\end{tcolorbox}


% =============================================================================
% PARAMETERS
% =============================================================================

\section{Parameter Calibration}
\label{sec:cps-parameters}

Table \ref{tab:cps-parameters} summarizes literature-validated parameters:

\begin{table}[htbp]
\centering
\caption{CPS Parameters with Literature Sources}
\label{tab:cps-parameters}
\begin{tabular}{lllll}
\toprule
\textbf{Parameter} & \textbf{Value} & \textbf{68\% CI} & \textbf{Source} \\
\midrule
$A_{\max}$ & 0.85 & [0.75, 0.95] & Sweller (1988) \\
$k$ (decay rate) & 1.5 & [1.2, 1.8] & LLMMC Prior \\
$\gamma_{\text{avg}}$ & +0.45 & [0.33, 0.57] & Funke (2001) \\
$\eta_{\text{start}}$ & 0.50 & [0.35, 0.65] & Dörner (1983) \\
$\eta_{\text{process}}$ & 0.60 & [0.45, 0.75] & Dörner (1983) \\
$\delta$ (delay) & 3--8 steps & Range & MicroDYN \\
$\theta_0$ & 0.55 & [0.45, 0.65] & Simon (1955) \\
$\tau$ (learning rate) & 0.15 & [0.10, 0.20] & LLMMC Prior \\
WMC → CPS ($\beta$) & 0.40 & [0.35, 0.45] & Greiff (2013) \\
$g$ → CPS ($\beta$) & 0.45 & [0.40, 0.55] & Greiff (2013) \\
\bottomrule
\end{tabular}
\end{table}


% =============================================================================
% SENSITIVITY ANALYSIS
% =============================================================================

\section{Sensitivity Analysis}
\label{sec:cps-sensitivity}

Parameter variation of ±20\% reveals the following influence on $U_{\text{CPS}}$:

\begin{table}[htbp]
\centering
\caption{Sensitivity Analysis: Influence on CPS Utility}
\label{tab:cps-sensitivity}
\begin{tabular}{llll}
\toprule
\textbf{Parameter} & \textbf{Influence} & \textbf{Rank} & \textbf{Interpretation} \\
\midrule
$\eta$ (Intransparency) & 42\% & 1 & \textbf{Main driver} \\
$\gamma$ (Connectivity) & 28\% & 2 & Second most important \\
$A_{\max}$ (Awareness) & 17\% & 3 & Moderate influence \\
$\delta$ (Delay) & 10\% & 4 & Lower influence \\
$\theta$ (Threshold) & 3\% & 5 & Minimal \\
\bottomrule
\end{tabular}
\end{table}

\textbf{Key Insight:} Intransparency (η) is the dominant lever. Interventions should prioritize \textbf{transparency enhancement} over other approaches.


% =============================================================================
% FRAMEWORK COMPARISON
% =============================================================================

\section{Framework Comparison: EBF vs. PISA 2015}
\label{sec:cps-comparison}

\begin{table}[htbp]
\centering
\caption{EBF 10C vs. PISA 2015 CPS Framework}
\label{tab:framework-comparison}
\begin{tabular}{lll}
\toprule
\textbf{Aspect} & \textbf{EBF 10C} & \textbf{PISA 2015 CPS} \\
\midrule
Focus & Utility-based (Why?) & Competency-based (What?) \\
Dimensions & 10 COREs & 12-cell matrix (4 × 3) \\
Social aspects & HOW ($\gamma$), WHO (L) & Explicit: Collaboration \\
Intransparency & $\eta$ (explicit) & Implicit in tasks \\
Strength & Explains mechanisms & Measures outcomes \\
Primary use & Intervention design & Assessment \\
\bottomrule
\end{tabular}
\end{table}

The PISA 2015 framework uses a 12-cell matrix crossing 4 cognitive processes (Exploring, Representing, Planning, Monitoring) with 3 social dimensions (Shared Understanding, Team Organisation, Taking Action). The EBF approach is complementary: it explains \textit{why} certain CPS behaviors occur, while PISA measures \textit{what} competencies are present.


% =============================================================================
% INTERVENTION GUIDELINES
% =============================================================================

\section{Intervention Guidelines}
\label{sec:cps-interventions}

Based on the sensitivity analysis, interventions should be prioritized as follows:

\begin{table}[htbp]
\centering
\caption{CPS Intervention Priorities (Ranked by Leverage)}
\label{tab:cps-interventions}
\begin{tabular}{cllll}
\toprule
\textbf{Rank} & \textbf{Target} & \textbf{Intervention Type} & \textbf{Examples} & \textbf{10C} \\
\midrule
1 & $\eta \downarrow$ & Transparency & Visualizations, Dashboards, Concept Maps & WHEN \\
2 & $\gamma \rightarrow$ & Connectivity & ``What affects what?'' training & HOW \\
3 & $A \uparrow$ & Load reduction & Chunking, external memory, time relief & AWARE \\
4 & $\delta \downarrow$ & Feedback speed & Simulation, rapid prototyping & WHEN \\
5 & $U_D \uparrow$ & Motivation & Mastery goals, success feedback & WHAT \\
\bottomrule
\end{tabular}
\end{table}


% =============================================================================
% WORKED EXAMPLES
% =============================================================================

\section{Worked Examples}
\label{sec:cps-examples}

\subsection{Example 1: Education (Complex Math Problems)}

\textbf{Problem:} Students fail at complex word problems.

\textbf{10C Diagnosis:}
\begin{itemize}
\item $\eta = 0.7$ (high intransparency of problem structure)
\item $A = 0.3$ (cognitive overload from text volume)
\item $\gamma = 0.5$ (connected variables not recognized)
\end{itemize}

\textbf{10C Intervention:}
\begin{enumerate}
\item $\eta \downarrow$: Concept maps visualizing variable relationships
\item $A \uparrow$: Chunking (decompose into sub-problems)
\item $\gamma \rightarrow$: Explicit training: ``What depends on what?''
\end{enumerate}

\textbf{Expected Effect:} $U_{\text{CPS}} \uparrow$ by approximately 35\% based on sensitivity weights.


\subsection{Example 2: Management (Strategic Decisions)}

\textbf{Problem:} Strategic errors in dynamic markets (cf. Dörner's ``Logic of Failure'')

\textbf{10C Diagnosis:}
\begin{itemize}
\item $\delta = 8+$ steps (long delay between decision and visible effect)
\item $\eta = 0.6$ (market dynamics not fully visible)
\item HIERARCHY: L1 decisions without L2 consistency check
\end{itemize}

\textbf{10C Intervention:}
\begin{enumerate}
\item $\delta \downarrow$: Leading indicators instead of lagging indicators
\item $\eta \downarrow$: Dashboards, scenario planning, war gaming
\item HIERARCHY: Implement Strategic Coherence Index (SCI)
\end{enumerate}


% =============================================================================
% SUMMARY BOX
% =============================================================================

\section{Summary}
\label{sec:cps-summary}

\begin{tcolorbox}[colback=green!5!white, colframe=green!75!black,
    title=Summary: CPS in the EBF]

\textbf{Main Result:} Complex Problem Solving is a \textbf{configuration} of 6 CORE dimensions (AWARE, HOW, WHEN, READY, STAGE, WHAT), not a separate theory.

\textbf{Core Formula:}
\begin{equation*}
U_{\text{CPS}}(t) = A(\Psi_{\text{load}}) \times (1-\eta) \times \left[U_D + \sum \gamma_{ij} \times \frac{dX_i}{dt} \times \frac{dX_j}{dt}\right]
\end{equation*}

\textbf{Key Insight:} Intransparency ($\eta$) is the main lever (42\%). Interventions should prioritize transparency enhancement.

\textbf{Intervention Priority:} $\eta \downarrow$ → $\gamma \rightarrow$ → $A \uparrow$ → $\delta \downarrow$ → $U_D \uparrow$

\end{tcolorbox}


% =============================================================================
% GLOSSARY
% =============================================================================

\section{Glossary of Symbols}
\label{sec:cps-glossary}

\begin{tabular}{ll}
\toprule
\textbf{Symbol} & \textbf{Definition} \\
\midrule
$A(\Psi_{\text{load}})$ & Awareness function under cognitive load \\
$A_{\max}$ & Maximum awareness capacity (≈ 0.85) \\
$\eta$ & Intransparency parameter $\in [0,1]$ \\
$\gamma_{ij}$ & Complementarity/connectivity between variables $i$ and $j$ \\
$\delta$ & Feedback delay (in steps or time units) \\
$\theta$ & Action threshold (satisficing level) \\
$\Psi_{\text{load}}$ & Cognitive load context dimension \\
$U_D$ & Development utility (competence motivation) \\
$U_{\text{CPS}}$ & Complex Problem Solving utility \\
WAX & Willingness to Act \\
\bottomrule
\end{tabular}

\vspace{1em}
For the complete glossary, see \textbf{Appendix G: Master Glossary}.


% =============================================================================
% CRITICAL FOUNDATIONS
% =============================================================================

\section{Critical Foundations and Potential Objections}
\label{sec:cps-foundations}

\subsection{Theoretical Foundations}

The 10C-CPS integration rests on the following foundational assumptions:

\begin{enumerate}
\item \textbf{Configuration Hypothesis:} CPS is not ontologically distinct from other decision-making processes but represents a specific parameter configuration within the universal 10C framework.

\item \textbf{Cognitive Load as Binding Constraint:} Working memory capacity limits awareness, making $A(\Psi_{\text{load}})$ the primary bottleneck in complex problem solving \citep{PAP-sweller1988cognitiveload}.

\item \textbf{Multiplicative Transparency:} Complete intransparency ($\eta = 1$) renders problem-solving impossible, justifying the multiplicative rather than additive structure.
\end{enumerate}

\subsection{Potential Objections}

\begin{description}
\item[Objection 1: Reductionism] \textit{``CPS has unique emergent properties that cannot be reduced to component dimensions.''}

\textbf{Response:} The 10C approach does not deny emergence but provides a generative mechanism. The specific configuration creates CPS-typical behaviors without requiring sui generis constructs.

\item[Objection 2: Domain Specificity] \textit{``CPS performance is domain-specific; a single model cannot capture this.''}

\textbf{Response:} Domain specificity enters through the Ψ-dimensions (context). The core mechanism remains constant; only parameters vary across domains.

\item[Objection 3: Social Neglect] \textit{``Individual CPS ignores collaborative aspects measured in PISA 2015.''}

\textbf{Response:} Social dimensions are addressed through WHO (levels L > 1) and HOW (social complementarity $\gamma_{social}$). A full collaborative CPS extension is identified as an open issue.
\end{description}


% =============================================================================
% OPEN ISSUES
% =============================================================================

\section{Open Issues and Research Agenda}
\label{sec:cps-open}

\begin{enumerate}
\item \textbf{Collaborative CPS:} The current model focuses on individual CPS. Integration with social CORE dimensions (WHO: L > 1) for collaborative problem solving requires further development.

\item \textbf{Domain Specificity:} Are CPS parameters ($\eta$, $\gamma$, $\delta$) domain-general or domain-specific? MicroDYN assumes domain-generality, but empirical evidence is mixed.

\item \textbf{Learning Dynamics:} How does $\eta$ decrease over time as problem-solvers learn the system? The current model is static; a dynamic version would specify $d\eta/dt$.

\item \textbf{Validation:} The parameter values are derived from LLMMC priors and literature review. Field validation in education and management contexts is needed.
\end{enumerate}


% =============================================================================
% REFERENCES
% =============================================================================

\section{References}
\label{sec:cps-references}

This appendix draws on 9 key papers in the CPS literature, all integrated into the master bibliography with full EBF 6-field classification:

\begin{itemize}[nosep]
\item \cite{PAP-doerner1989logik} -- Foundational CPS theory
\item \cite{PAP-doerner1983lohhausen} -- Original Lohhausen study
\item \cite{PAP-funke2017cps} -- CPS definition and review
\item \cite{PAP-greiff2013cps} -- MicroDYN validation (N=855)
\item \cite{PAP-greiff2012microdynamvalidity} -- Incremental validity
\item \cite{PAP-greiff2017microdynamapproach} -- MicroDYN methodology
\item \cite{PAP-oecd2017problemsolving} -- OECD problem solving report
\item \cite{PAP-oecd2017pisa2015cps} -- PISA 2015 CPS assessment
\item \cite{PAP-sweller1988cognitiveload} -- Cognitive Load Theory
\end{itemize}

\vspace{1em}
\noindent For additional references, see \textbf{bcm\_master.bib} (Master Bibliography).

\nocite{bcm_master}


% =============================================================================
% CROSS-REFERENCE FOOTER
% =============================================================================

\vspace{2em}
\hrule
\vspace{1em}

\textbf{Cross-References:}
\begin{itemize}[nosep]
\item \textbf{CORE Appendices:} AU (AWARE), V (WHEN), B (HOW), C (WHAT), AV (READY), AW (STAGE), HI (HIERARCHY)
\item \textbf{Method Appendices:} HHH (METHOD-TOOLKIT), CAL (METHOD-LLMMC)
\item \textbf{Related Chapters:} 11 (AWARE), 9 (WHEN), 4 (HOW), 15 (HIERARCHY)
\item \textbf{Model Registry:} MOD-CPS-001
\item \textbf{Session:} EBF-S-2026-01-28-EDU-001
\end{itemize}

\end{document}
